%==============================================================
%  Currículo — Hugo Guimarães da Silva
%  Compilar com pdfLaTeX (Overleaf: Menu > Compiler > pdfLaTeX)
%  Tema verde para combinar com o portfólio.
%==============================================================
\documentclass[a4paper,10pt]{article}

\usepackage[T1]{fontenc}
\usepackage[utf8]{inputenc}
\usepackage[brazil]{babel}
\usepackage[default]{lato}     % fonte sans moderna (se não compilar, comente esta linha)
\usepackage[top=1.1cm,bottom=1.1cm,left=1.4cm,right=1.4cm]{geometry}
\usepackage{xcolor}
\usepackage{titlesec}
\usepackage{enumitem}
\usepackage{fontawesome5}
\usepackage[hidelinks]{hyperref}
\usepackage{tabularx}

%--------------------------------------------------------------
%  Cores e estilo
%--------------------------------------------------------------
\definecolor{accent}{HTML}{16A34A}   % verde
\definecolor{dark}{HTML}{1A1A1A}
\definecolor{muted}{HTML}{5A5A5A}

\hypersetup{colorlinks=true, urlcolor=accent, linkcolor=accent}

% Cabeçalho de seção com linha colorida
\titleformat{\section}
  {\large\bfseries\color{accent}}
  {}{0em}{}
\titlespacing*{\section}{0pt}{8pt}{3pt}

\setlist[itemize]{leftmargin=1.2em, itemsep=1pt, topsep=1pt, parsep=0pt}
\pagestyle{empty}
\setlength{\parindent}{0pt}
\setlength{\parskip}{0pt}

%--------------------------------------------------------------
%  Comandos auxiliares
%--------------------------------------------------------------
% Entrada: {cargo/título}{período}{instituição/subtítulo}
\newcommand{\entry}[3]{%
  \textbf{#1}\hfill{\color{muted}\small #2}\\
  {\color{muted}\itshape #3}\par\vspace{2pt}
}

% Projeto: {nome}{tag}{stack}{descrição}{link-texto}{link-url}
\newcommand{\project}[6]{%
  \textbf{#1} \hspace{4pt}{\footnotesize\color{accent}\faTag\ #2}\\
  {\color{muted}\footnotesize #3}\par\vspace{1pt}
  #4\ \href{#6}{\footnotesize\color{accent}\faExternalLinkAlt\ #5}\par\vspace{2pt}
}

%==============================================================
\begin{document}

%--------------------------------------------------------------
%  Cabeçalho
%--------------------------------------------------------------
{\Huge\bfseries\color{dark} Hugo Guimarães da Silva}\\[3pt]
{\large\color{accent} Desenvolvedor Full Stack}\\[6pt]

\small
\faMapMarker*\ Paranaguá, PR, Brasil \quad
\faPhone*\ (41) 98502-5645 \quad
\faEnvelope\ \href{mailto:hugoartess@gmail.com}{hugoartess@gmail.com}\\[2pt]
\faLinkedin\ \href{https://linkedin.com/in/hugoGsilvaa}{linkedin.com/in/hugoGsilvaa} \quad
\faGithub\ \href{https://github.com/HugoGsilva}{github.com/HugoGsilva} \quad
\faGlobe\ \href{https://hugogsilva.dev}{hugogsilva.dev}
\normalsize

%--------------------------------------------------------------
%  Resumo
%--------------------------------------------------------------
\section{Resumo}
Desenvolvedor full stack e estudante de Análise e Desenvolvimento de Sistemas. Construo
APIs, automações e aplicações web de ponta a ponta, da interface ao deploy, com foco em
código testável e em produção. Tenho experiência com infraestrutura própria (Docker,
Portainer, Cloudflare) e, atualmente, exploro IA aplicada e automações com n8n.

%--------------------------------------------------------------
%  Experiência
%--------------------------------------------------------------
\section{Experiência}
\entry{Desenvolvedor Full Stack}{Desde 2026}{Autônomo / Projetos próprios}
\begin{itemize}
  \item Desenvolvimento de aplicações web, APIs REST e automações, com entrega completa: frontend, backend, persistência, integração e deploy.
  \item Aplicação em produção real para automação de propostas comerciais com geração de PDF.
\end{itemize}
\vspace{2pt}
\entry{Infraestrutura e DevOps (estudos e projetos pessoais)}{Desde 2025}{Autônomo}
\begin{itemize}
  \item Deploy e administração de aplicações com Docker e Portainer em VPS Linux.
  \item DNS, SSL e proxy reverso via Cloudflare e Nginx; automações de fluxos com n8n.
\end{itemize}

%--------------------------------------------------------------
%  Projetos
%--------------------------------------------------------------
\section{Projetos em Destaque}
\project{Gerador de Propostas Comerciais}{Produção}
  {React · Node.js · Express · Docxtemplater · LibreOffice · Docker}
  {Aplicação full stack que preenche modelos DOCX/ODT e gera PDFs comerciais, automatizando um processo manual e repetitivo. Em uso real, com deploy próprio.}
  {docx.hugogsilva.dev}{https://docx.hugogsilva.dev}

\project{IndieGameStore API}{Acadêmico}
  {Ruby on Rails 8 · PostgreSQL · JWT/RBAC · Swagger · React · Vite}
  {Marketplace de jogos indie com autenticação JWT, controle de acesso por papéis, reviews, compras, multi-moeda e documentação Swagger. Frontend React.}
  {loja.hugogsilva.dev}{https://loja.hugogsilva.dev}

\project{Chatbot Netflix-Prolog}{Acadêmico}
  {FastAPI · SWI-Prolog · Redis · Angular 17 · Docker Compose · Cypress}
  {Chatbot com NLU no servidor (correção ortográfica e fuzzy matching), inferência simbólica em Prolog, sessões e rate limiting em Redis, frontend Angular.}
  {chat.hugogsilva.dev}{https://chat.hugogsilva.dev}

%--------------------------------------------------------------
%  Formação
%--------------------------------------------------------------
\section{Formação}
\entry{Análise e Desenvolvimento de Sistemas}{Em andamento}{Instituto Federal do Paraná (IFPR), Paranaguá, PR}
\vspace{2pt}
\entry{Técnico em Comércio Exterior}{Concluído em 2022}{Colégio Alberto Gomes Veiga, Paranaguá, PR}

%--------------------------------------------------------------
%  Habilidades
%--------------------------------------------------------------
\section{Habilidades}
\begin{tabularx}{\textwidth}{@{}l X@{}}
  \textbf{Frontend}      & HTML, CSS, JavaScript, React, Angular, Vite, Tailwind CSS \\[3pt]
  \textbf{Backend}       & Node.js, Express, Python, FastAPI, Ruby on Rails, Java \\[3pt]
  \textbf{Dados}         & SQL, PostgreSQL, Redis, REST, Swagger \\[3pt]
  \textbf{Infra \& Deploy} & Docker, Portainer, Cloudflare, Nginx, Linux/VPS, Git, CI/CD \\[3pt]
  \textbf{Automação}     & n8n, AutoHotkey, scripts utilitários, geração de documentos \\[3pt]
  \textbf{Explorando}    & Treinamento/fine-tuning de modelos (noções), APIs de LLM, Prompt engineering \\
\end{tabularx}

%--------------------------------------------------------------
%  Idiomas
%--------------------------------------------------------------
\section{Idiomas}
\textbf{Português:} Nativo \hspace{2em} \textbf{Inglês:} Avançado

\end{document}
